\documentclass[a4paper,fleqn]{cas-sc}
\usepackage[numbers]{natbib}
\usepackage{amsmath,amssymb,amsthm}
\newtheorem{cbmptheorem}{Định lý}
\usepackage{graphicx}
\usepackage{iftex}
\ifPDFTeX
  \usepackage[T5]{fontenc}
\else
  \usepackage[no-math]{fontspec}
  \setmainfont{texgyretermes-regular.otf}[BoldFont=texgyretermes-bold.otf,
    ItalicFont=texgyretermes-italic.otf,BoldItalicFont=texgyretermes-bolditalic.otf]
\fi
\graphicspath{{figures/}{../results/figures/}{../results/summary/}}
\begin{document}
\shorttitle{Exact methods for cyclic bandwidth}
\shortauthors{CBMP project team 2}
\title[mode=title]{SAT Encodings for Cyclic Bandwidth Minimization}
\author{Nguyen Minh Hung - 24021496}
\author{Nguyen Minh Duc - 24021416}
\author{Ninh Ngoc Hung - 24021504}
\begin{abstract}
We study the minimum possible maximum edge distance when the vertices of
an undirected graph occupy distinct positions on a circle. We formulate
the bounded-distance decision problem in SAT and use a decreasing sequence
of thresholds to prove optimality or return a validated layout and lower
bound at timeout. A pilot on three graphs from Harwell--Boeing matrices
illustrates both outcomes. For the 48-vertex graph derived from BCSSTK01,
the pairwise solver and CP-SAT each found bandwidth 12 within the 60-second
budget; the recorded lower bound was 6.
\end{abstract}
\begin{keywords}
Cyclic bandwidth \sep SAT encoding \sep graph labeling \sep constraint programming
\end{keywords}
\maketitle

\section{Introduction}
% Motivate the problem and state the research questions.

\section{Background}
A layout assigns the $n$ vertices of a simple undirected graph $G=(V,E)$
to distinct positions $0,\ldots,n-1$ on a circle. For positions $i,j$ the
distance is $d_n(i,j)=\min(|i-j|,n-|i-j|)$. The bandwidth of a layout
is the largest distance across an edge, with value zero for an empty edge
set. The optimization problem minimizes this bandwidth over all layouts.

For an integer threshold $K$, the decision problem asks whether a layout
of bandwidth at most $K$ exists. Each position has at most $2K$ distinct
neighbours within distance $K$, so $\lceil\Delta(G)/2\rceil$ is a lower
bound on the optimum. Every permutation has bandwidth at most
$\lfloor n/2\rfloor$. A validated layout of value $B$ gives an upper bound.
Equality with a proven lower bound, or a proof of infeasibility at $B-1$,
certifies optimality. An interrupted decision leaves that threshold unresolved.

\section{Pairwise encoding}
% Describe variables, CNF clauses and correctness.

\section{Order encoding}
% Describe threshold variables, clauses and decoding.
\subsection{Chứng minh tính đúng đắn của mã hóa Order}

Xét đồ thị vô hướng $G = (V, E)$ với $|V| = n$ đỉnh (đánh số $0, 1, \dots, n-1$), tập cạnh $E$ và ngưỡng băng thông vòng nguyên $K$. Ký hiệu khoảng cách vòng giữa hai vị trí $i, j \in \{0, 1, \dots, n-1\}$ là
\[
d_c(i, j) = \min\big(|i - j|,\; n - |i - j|\big).
\]
Trong \texttt{order.py}, với mỗi đỉnh $v \in \{0, \dots, n-1\}$ và ngưỡng $t \in \{0, \dots, n-2\}$, biến mệnh đề dương \texttt{y[v, t]} (tương ứng chỉ số \texttt{threshold(v, t, n)}) mang ý nghĩa:
\[
\texttt{y[v, t]} \iff \big(p(v) \le t\big).
\]
Gọi $F(K)$ là tập mệnh đề CNF sinh bởi \texttt{build(graph, K)} (với $K \ge 0$ và $n \ge 2$). Ta chứng minh rằng $F(K)$ là \textbf{SAT} khi và chỉ khi tồn tại một hoán vị (cách xếp) $p : V \to \{0, 1, \dots, n-1\}$ có cyclic bandwidth không vượt quá $K$, tức $\max_{(u,v) \in E} d_c(p(u), p(v)) \le K$.

\subsubsection{Chiều thuận: Từ cách xếp sang SAT}

Giả sử tồn tại một hoán vị hợp lệ $p$ trên $\{0, 1, \dots, n-1\}$ sao cho $d_c(p(u), p(v)) \le K$ với mọi cạnh $(u, v) \in E$. Ta thiết lập phép gán chân lý cho các biến của \texttt{order.py} bằng quy tắc:
\[
\texttt{y[v, t]} = \text{True} \iff p(v) \le t, \quad \forall v \in \{0, \dots, n-1\},\; \forall t \in \{0, \dots, n-2\}.
\]
Phép gán này thỏa mãn toàn bộ mệnh đề trong $F(K)$ vì các lý do sau:
\begin{itemize}
    \item \textbf{Mệnh đề đơn điệu (\texttt{[-var\_map[v, t], var\_map[v, t + 1]]}):} Nếu $p(v) \le t$ thì hiển nhiên $p(v) \le t + 1$, nghĩa là \texttt{y[v, t]} đúng kéo theo \texttt{y[v, t + 1]} đúng. Do đó mọi mệnh đề $(\neg \texttt{y[v, t]} \lor \texttt{y[v, t + 1]})$ đều được thỏa mãn.
    \item \textbf{Mệnh đề khác vị trí (\texttt{different\_from}):} Với mỗi đỉnh $v$ và vị trí $i \in \{0, \dots, n-1\}$, \texttt{different\_from(v, i, n)} biểu diễn điều kiện $p(v) \neq i$ thông qua các ngưỡng $(p(v) \le i-1)$ khi $i > 0$ và $\neg(p(v) \le i)$ khi $i < n-1$. Vì $p$ là một hoán vị nên với mọi cặp đỉnh phân biệt $(u, v)$ và mọi vị trí $i$, ta không thể có đồng thời $p(u) = i$ và $p(v) = i$; ít nhất một trong hai điều kiện $p(u) \neq i$ hoặc $p(v) \neq i$ đúng, nên mệnh đề \texttt{different\_from(u, i, n) + different\_from(v, i, n)} luôn đúng.
    \item \textbf{Mệnh đề khoảng cách vòng $\le K$:} Nếu $K \ge \lfloor n/2 \rfloor$, \texttt{build} không phát sinh mệnh đề cấm cặp vị trí nào và mọi $d_c(p(u), p(v)) \le \lfloor n/2 \rfloor \le K$. Nếu $K < \lfloor n/2 \rfloor$, với mọi cạnh $(u, v) \in E$ ta có $d_c(p(u), p(v)) \le K$. Do đó cặp vị trí thực tế $(i, j) = (p(u), p(v))$ không bao giờ rơi vào trường hợp $d_c(i, j) > K$, tức là không vi phạm bất kỳ mệnh đề cấm \texttt{different\_from(u, i, n) + different\_from(v, j, n)} nào.
\end{itemize}
Vậy $F(K)$ là \textbf{SAT}.

\subsubsection{Chiều đảo: Từ SAT về cách xếp}

Ngược lại, giả sử $F(K)$ là \textbf{SAT} với một mô hình thỏa mãn $\mathcal{M}$. Hàm \texttt{decode} lấy ngưỡng $t \in \{0, \dots, n-2\}$ nhỏ nhất có \texttt{y[v, t]} đúng (trong tập \texttt{positive}), hoặc trả về $n - 1$ nếu toàn bộ chuỗi ngưỡng của $v$ đều sai.

\begin{itemize}
    \item \textbf{Mỗi chuỗi ngưỡng giải mã ra đúng một vị trí (kể cả hai trường hợp biên khi $n \ge 2$):} Do các mệnh đề đơn điệu $(\neg \texttt{y[v, t]} \lor \texttt{y[v, t + 1]})$ ép các giá trị đúng nằm ở một hậu tố liên tiếp trên dãy ngưỡng $t = 0, 1, \dots, n-2$, nên tồn tại duy nhất một giá trị giải mã $p(v) \in \{0, 1, \dots, n-1\}$ tương thích với dãy ngưỡng:
    \begin{itemize}
        \item \emph{Trường hợp biên chuỗi toàn đúng (vị trí $0$):} Nếu \texttt{y[v, 0]} đã đúng thì do tính đơn điệu, mọi \texttt{y[v, t]} với $t \in \{0, \dots, n-2\}$ đều đúng; \texttt{decode} chọn ngưỡng nhỏ nhất là $0$, ứng với $p(v) = 0$.
        \item \emph{Trường hợp biên chuỗi toàn sai (vị trí $n-1$):} Khi $n \ge 2$, nếu mọi \texttt{y[v, t]} với $t \in \{0, \dots, n-2\}$ đều sai (tập \texttt{positive} không chứa ngưỡng nào của $v$), \texttt{decode} trả về giá trị mặc định $n - 1$, ứng với $p(v) = n - 1$ (phù hợp với quy ước ngầm định $p(v) \le n-1$ luôn đúng).
    \end{itemize}
    Với mọi giá trị trung gian $p(v) = t^* \in \{1, \dots, n-2\}$, ta có \texttt{y[v, t]} sai với mọi $t < t^*$ và đúng với mọi $t \ge t^*$. Như vậy, mỗi đỉnh $v$ luôn nhận duy nhất một vị trí $p(v) \in \{0, \dots, n-1\}$.

    \item \textbf{Các đỉnh nhận vị trí khác nhau (hoán vị):} Giả sử tồn tại hai đỉnh phân biệt $u \neq v$ cùng được giải mã ra cùng một vị trí $i \in \{0, \dots, n-1\}$. Khi đó cả \texttt{different\_from(u, i, n)} và \texttt{different\_from(v, i, n)} đều có mọi literal đánh giá là sai dưới mô hình $\mathcal{M}$ (vì với $i > 0$ thì \texttt{y[u, i-1]} sai do $p(u) = i > i-1$, và với $i < n-1$ thì \texttt{y[u, i]} đúng do $p(u) = i \le i$ nên \texttt{-y[u, i]} sai; tương tự cho $v$). Điều này mâu thuẫn với mệnh đề \texttt{different\_from(u, i, n) + different\_from(v, i, n)} trong $F(K)$. Vì có $n$ đỉnh nhận giá trị trong tập $n$ phần tử $\{0, \dots, n-1\}$ và đôi một khác nhau, $p$ là một hoán vị hợp lệ.

    \item \textbf{Mọi cạnh đều có khoảng cách vòng $\le K$:} Với $K < \lfloor n/2 \rfloor$, giả sử tồn tại cạnh $(u, v) \in E$ sao cho $d_c(p(u), p(v)) > K$. Đặt $i = p(u)$ và $j = p(v)$. Khi đó mệnh đề \texttt{different\_from(u, i, n) + different\_from(v, j, n)} có mặt trong $F(K)$, nhưng mọi literal của cả hai vế đều sai dưới $\mathcal{M}$ vì $p(u) = i$ và $p(v) = j$ — mâu thuẫn với việc $\mathcal{M}$ thỏa mãn $F(K)$. Với $K \ge \lfloor n/2 \rfloor$, mọi cặp vị trí trên đồ thị $n$ đỉnh đều có khoảng cách vòng $\le \lfloor n/2 \rfloor \le K$. Do đó $\max_{(u,v) \in E} d_c(p(u), p(v)) \le K$.
\end{itemize}

\subsubsection{Kết luận}

Từ hai chiều trên, với mọi ngưỡng $K \ge 0$ (và $n \ge 2$), công thức $F(K)$ sinh bởi \texttt{order.py} là \textbf{SAT} khi và chỉ khi tồn tại một cách xếp (hoán vị) $p$ các đỉnh của đồ thị sao cho cyclic bandwidth của $p$ không vượt quá $K$.
\section{Baselines}
% Document CP and MIP formulations and implementations.

\section{Experiments}
\label{sec:experiments}

\subsection{Benchmark graphs and preprocessing}
The selected benchmark graphs are drawn from the Harwell--Boeing collection.
The raw matrix files and normalized graph files are stored separately. Each
entry in the versioned manifest records the source page, raw-file path, graph
path, matrix type, number of vertices, number of normalized undirected edges,
maximum degree, number of isolated vertices, and SHA-256 checksums for both
raw and normalized data. A further checksum is computed from the canonical
sequence consisting of $n$ and the sorted edge list.

For a square matrix $A$, vertices correspond to rows and columns. An undirected
edge $\{i,j\}$ is included for $i\ne j$ if at least one of $A_{ij}$ or
$A_{ji}$ is nonzero. Diagonal entries are discarded, duplicate and
opposite-direction entries are merged, and isolated vertices are retained.
Consequently, the number of nonzero matrix entries is not generally the number
of edges in the normalized simple graph. For example, the source reports 50
stored matrix entries for JGL009~\cite{NISTJGL009}, whereas the normalized graph used here has 32 unique
undirected edges. We therefore do not copy the paper's matrix-entry counts into
the normalized manifest.

\begin{table}[t]
\centering
\caption{Three Harwell--Boeing graphs used in the 300-second main benchmark.}
\label{tab:dataset}
\begin{tabular}{lrrrrr}
\hline
Instance & $n$ & $m$ & $\Delta(G)$ & Isolates & Density \\
\hline
JGL009 & 9 & 32 & 8 & 0 & 0.8889 \\
BCSSTK01 & 48 & 176 & 11 & 0 & 0.1560 \\
BP 400 & 822 & 4015 & 294 & 0 & 0.0119 \\
\hline
\end{tabular}
\end{table}

JGL009 and BCSSTK01 also occur in Table A.2 of Ren et al.~\cite{Ren2020CBMP}
and their source records are available from NIST~\cite{NISTJGL009,NISTBCSSTK01}.
BP 400 is an additional Harwell--Boeing matrix from the SMTAPE set~\cite{NISTBP400}
and is used as a larger sparse instance. BCSSTM01 is documented separately as
an excluded case: after removing its diagonal entries, it has 48 isolated
vertices and no edges~\cite{NISTBCSSTM01}. The two distributed forms of JGL009
(plain and gzip-compressed) are the same matrix and are counted once. The
repository also records the 28 names in
Table A.2 as candidates; a candidate is included in a benchmark only
after its raw file, normalized graph, and checksums have been verified.

\subsection{Pilot protocol and observed outcomes}
The day-8 pilot covers three verified graphs and the five methods available in the original run: pairwise SAT, order SAT, CP-SAT, CPLEX MIP, and
Gurobi MIP. Each planned run has a 60-second end-to-end wall-clock budget,
one thread, and seed zero. The original run log contains 15 records. Table~\ref{tab:day8-pilot}
uses the 60-second CPLEX Studio 22.2 reruns on JGL009 and BCSSTK01
and the 60-second Gurobi 13.0.3 rerun on all three graphs. In the table, OPT means
optimality proven, F/TO means a validated feasible incumbent with timeout, and
Guard means rejected by the preflight size guard before solver startup.

On JGL009 the identity bandwidth and the degree lower bound both equal 4. Both SAT runs stopped before CNF construction. On BCSSTK01,
OR-Tools CP-SAT proved bandwidth 12 optimal in 46.686 s (record 8 of \path{results/pilot_day8.jsonl}). In the earlier day-6 pilot, pairwise SAT and CP-SAT each reached 12 with lower bound 6. Both SAT encodings return
validated feasible solutions but do not prove optimality within 60 seconds.
On BP 400, CP-SAT returns a feasible solution but does not prove optimality
within the time limit; the pairwise and order SAT encodings are rejected by
the preflight guard because their estimated clause counts are approximately
3.27 billion and 2.99 billion, respectively. The guard adds $mn^2$ possible edge clauses to each assignment count before choosing $K$. This upper estimate includes every ordered position pair; solver time for a guarded run is N/A.

The new 60-second CPLEX Studio reruns proved bandwidth 4 optimal on
JGL009 and returned a validated feasible layout of bandwidth 24 on
BCSSTK01, with lower bound 6 and timeout. Thus the latter entry has
bounds $[6,24]$ and carries no optimality proof. The Gurobi academic-license rerun proved bandwidth 4 optimal
on JGL009. On BCSSTK01 it returned bounds $[6,17]$ at timeout. BP~400 exceeded
the runner's size guard before either MIP solver started. Both MIP solvers
also returned validated layouts on the smaller graphs under the separate
300-second protocol (Table~\ref{tab:six-method-rerun}).

The five-method pilot excludes IBM CP Optimizer. The separate six-method
60-second pilot applies a 1 GiB Windows job memory cap. Gurobi exhausted the memory allowance on BCSSTK01 before returning a layout. The independent
order/CP pilot is identified in the Order Encoding section.

\subsection{Main benchmark protocol}
The six-method rerun used an AMD Ryzen 5 5600H with 15.4 GiB of
OS-visible RAM, Windows 11 Home (build 26300), and Python 3.12.3.
Each instance--method pair had a 300-second target including model construction,
one requested solver thread, and seed zero where supported. Cleanup and
native solver return can add elapsed time beyond the target. The six-method
60-second pilot limited total Windows job committed memory to 1024 MiB
for the runner and its children. The 300-second run used model-size guards
without that memory cap. Some runs overlapped other local work; elapsed
times here describe these runs and are unsuitable for a speed ranking.

The benchmark runner stores one JSON object per run in JSONL format and
supports resuming completed runs. It preserves failures and distinguishes
solver status (optimal, feasible, or no solution) from termination reason
(completed, timeout, out-of-memory, license error, or other error). Every
reported incumbent is independently validated. The five-method, 60-second pilot and the six-method, 300-second run are
separate experiments. The pilot outcomes are reported above; the six-method
batch is summarized in the section on MIP and six-method results.

% Results and implementation analysis owned by Ninh Ngoc Hung.
% The historical 60-second pilot is described in experiments.tex;
% this section reports the shared MIP formulation and the six-method run.

\subsection{MIP model size and scalability}
The CPLEX MIP and Gurobi MIP adapters use the same assignment-based
formulation. For a graph with $n$ vertices and $m$ undirected edges, the
implementation creates $n^2+1$ variables ($n^2$ binary assignment variables
and one integer objective variable) and $2n+mn^2$ constraints. The latter
count follows the implementation's edge--position generation loop, including
rows for position pairs whose cyclic distance is zero.

\begin{table}[t]
\centering
\small
\caption{Estimated size of the shared MIP formulation on the verified graph
set. Counts are computed before invoking either optimizer.}
\label{tab:mip-model-size}
\begin{tabular}{lrrrr}
\hline
Instance & $n$ & $m$ & Variables & Constraints \\
\hline
JGL009 & 9 & 32 & 82 & 2,610 \\
BCSSTK01 & 48 & 176 & 2,305 & 405,600 \\
BP 400 & 822 & 4,015 & 675,685 & 2,712,872,904 \\
\hline
\end{tabular}
\end{table}

The $n^2$ assignment variables and $mn^2$ edge--position constraints make
this direct formulation impractical on BP 400. The runner's one-million-
constraint preflight limit rejects that instance before allocating the MIP
model. This is a configured resource guard, not an optimizer result, timeout,
or infeasibility proof.

\subsection{Six-method benchmark outcomes}
The consolidated table contains 18 method--instance entries. Fourteen
have validated layouts, including the two CPLEX Studio reruns. The
remaining four were stopped by the BP 400 size guards.
Table~\ref{tab:six-method-rerun} uses a 300-second budget and one thread
per entry. Sixteen entries come from the original batch at commit
\texttt{9df7be6}. The JGL009 and BCSSTK01 entries for IBM CPLEX MIP
come from Studio 22.2 reruns with the same formulation and budget.

All six methods proved 4 on JGL009. On BCSSTK01, pairwise SAT and both CP
engines reached 12, order SAT reached 13, Gurobi MIP reached 15, and CPLEX
MIP reached 18. All six runs ended with a reported lower bound of 6 and
without an optimality proof. CPLEX reported a solver bound of 1, which
the independent degree bound strengthened to 6. On BP 400, pairwise SAT and
order SAT were rejected by their estimated-size guards, and both MIP
adapters were rejected by the one-million-constraint guard before solver
startup. CP-SAT and CP Optimizer returned validated incumbents of 399 and
263, respectively, with certified lower bound 147, but neither proved
optimality within the time limit.

Solver runtime is N/A for the four guarded BP~400 runs.

\subsection{Experimental limitations}
The historical main manifest contains three graphs, not the full 28-instance
candidate list in Table A.2 of Ren et al. The set provides a small dense,
medium sparse, and larger sparse example, but is too small for broad
statistical or scalability claims. BP 400 is an additional Harwell--Boeing /
SMTAPE instance, not one of the 28 Table A.2 candidates. The graph edge count
$m$ is the number of unique undirected edges after removing self-loops and
merging duplicate or opposite-direction entries; it is not the raw matrix
nonzero count.

Both MIP solvers produced validated layouts on four matrix graphs:
JGL009, BCSSTK01, IBM32 and BCSPWR01. Only JGL009 yielded an optimality
proof from both. The other three comparisons measure the layouts and
bounds reached within each budget. BP~400 provides a model-size estimate
only, and the separate BCSPWR02 experiment evaluated the SAT encodings.


\subsection{Six-method rerun}
Table~\ref{tab:six-method-rerun} uses the three verified graphs and the
300-second run records. Every layout was rechecked against its graph
checksum. Reported lower bounds include the independent degree bound.
The original five-method experiment and the rerun used different solver
environments and settings; their proof statuses are reported separately.
On JGL009, Gurobi and CP Optimizer both certified 4 in the rerun. On
BCSSTK01, both CP engines reached bandwidth 12. The memory-limited
60-second pilot recorded a Gurobi allocation failure on this graph,
while the 300-second run without that cap returned a feasible layout.
\begin{table}[htbp]
\centering\small
\caption{Six-method rerun, 300 seconds per graph and method. Three graphs, one thread. Single values are proven optima. Brackets give the larger of the recorded and degree lower bounds and the validated incumbent at timeout. Size guard means rejection before solving (runtime N/A). The two smaller CPLEX MIP entries are separate Studio reruns.}
\label{tab:six-method-rerun}
\begin{tabular}{lccc}
\hline
Method & JGL009 & BCSSTK01 & BP 400 \\
\hline
Pairwise SAT & 4 & [6, 12] & Size guard \\
Order SAT & 4 & [6, 13] & Size guard \\
CP-SAT & 4 & [6, 12] & [147, 399] \\
CPLEX MIP & 4 & [6, 18] & Size guard \\
Gurobi MIP & 4 & [6, 15] & Size guard \\
CP Optimizer & 4 & [6, 12] & [147, 263] \\
\hline
\end{tabular}
\end{table}

\section{Conclusion}
Six solver adapters now share a benchmark runner with independent validation.
All six obtained bandwidth 1 on both toy graphs. The original five-method
benchmark certified bandwidth 12 on BCSSTK01 with CP-SAT. A later run on
the same graph may return a feasible layout of value 12 while leaving a
gap between its own bounds; each table therefore reports proof status per run.
On the scrambled eight-vertex path, reusing the order solver reduced the
total input clauses from 1320 to 552 over the same three thresholds.
The three-matrix dataset and the observed license, memory, and model-size
limits restrict conclusions about relative performance.


\section*{Data availability}
The project uses the graph inputs and conversion rules in the project
repository. The original pilot results are in \path{results/summary/day6_pairwise_cpsat.csv}.
The six-method run records and controlled SAT experiments are archived in
\path{results/raw/}.
\url{https://github.com/NguyenHung2604/BTL-2}

\bibliographystyle{cas-model2-names}
\bibliography{references}
\end{document}
